\section{Related Work}

\subsection{LLM-Guided Evolutionary Coding Approaches}
LLMs can act as mutation operators inside evolutionary loops over executable programs.
FunSearch \citep{romera-paredes2024} and AlphaEvolve \citep{novikov2025c,georgiev2025c} established this paradigm, and a growing collection of open-source frameworks now extend it, including OpenEvolve \citep{openevolve}, GEPA \citep{agrawal2026}, ShinkaEvolve \citep{lange2025b}, GigaEvo \citep{khrulkov2025}, CodeEvolve \citep{assumpcao2026}, the FM Agent \citep{li2026}, and AIDE \citep{jiang2025a}.
A second wave targets the search procedure itself by adapting strategies, models, or signals during the run \citep{liu2026,cemri2026,yan2026,ray2026,chen2026b}, or by changing the unit of evolution to solution spaces, strategies, skills, prompt groups, populations, or concept trees \citep{zhai2025,luo2026,ye2026,li2025c,zhang2025d,leleu2026,pourcel2026,wang2025a,yuksekgonul2026,ma2025,singhal2026,liu2026selfplayevolvesselfsyntheticpipeline}.

A particularly active line of work targets GPU kernel optimization, where wall-clock runtime provides a tight reward signal \citep{guo2025,wiedemann2026,cao2026,du2026,nichols2026,wei2025a,su2026,cheng2025a,cheng2025b}.
These frameworks have also been applied to compiler heuristics, computer architecture, cosmology, swarm-intelligence design, symbolic regression, retrieval, recommendation, hyperparameter optimization, code optimization, agentic reasoning, autonomous data science, and broader scientific discovery \citep{chen2026a,gupta2026,li2026a,cen2025,song2025,nian2026,wang2026a,ferreira2026,hu2026,zhou2026b,yang2025a,weng2026,skydiscover2026}.
This line builds on classical evolutionary computation and quality-diversity search \citep{koza1994,jaderberg2017,qian2024,nadizar2025,hughes2024}; because the evolved objects are programs, recent work develops program-aware notions of diversity and similarity \citep{friedman2023,zhang2026c}.
Our work is complementary: rather than proposing another framework, we analyze traces from four of them (OpenEvolve, GEPA, ShinkaEvolve, EvoX) to study the mechanisms by which programs improve, stagnate, diversify, or collapse.

\subsection{Analyzing Evolutionary Coding Agents}
A separate body of work studies evolutionary coding systems themselves.
Surveys situate the broader space \citep{fang2025,xie2025b}, benchmarks evaluate agents across ML competitions, long-horizon algorithm engineering, frontier research science, and production deployments \citep{chan2025b,imajuku2025,lupidi2026,ye2026a,pan2026,zheng2026a}, and \citet{gideoni2026} show that simple baselines can match elaborate evolutionary pipelines.
Closer to our methodology, several papers analyze search behavior rather than only final scores: trajectory analyses \citep{zhang2026d}, fitness-landscape characterization \citep{liu2025a}, failure modes of iterative LLM optimization \citep{nie2026}, exploration deficits \citep{pan2025}, output homogeneity \citep{jiang2025}, emergent risks in self-evolving agents \citep{shao2026}, and a taxonomy of multi-agent failures \citep{cemri2025}.
Our work adopts a similar diagnostic perspective but focuses on full evolutionary coding traces, which we use to characterize how populations evolve, how improvement propagates through lineages, and how search dynamics produce both successes and failure modes.
